% GENERATED FILE. Edit canonical lessons, then run npm run generate:books.
% canonical-source-hash: fnv1a64:f96220ab9970120f
% canonical-lessons: SA-C555-prerana, SA-C555-udyuktah, SA-C555-pratyuta, SA-C555-karmathah, SA-C555-adhyavasyati, SA-R555-review-549-555, SA-R555-review-final-555

\chapter{The Road Ahead}
\label{ch:sa-the-road-ahead}
\hlchaptermodality{555}

\begin{chapteropening}
\textbf{By the end of this chapter:} \emph{I can say what I resolve to do and what will follow.}

Complete the last lesson of chapter 555: I can say what I resolve to do and what will follow.
\end{chapteropening}

\section[\texorpdfstring{\sk{प्रेरणा} (preraṇā)}{प्रेरणा (preraṇā)}]{\sk{प्रेरणा} (preraṇā) — inspiration; urging}

\label{lesson:SA-C555-prerana}



\begin{quote}
Before the new one: say the Sanskrit for the time ahead, then the Sanskrit for a start.
\end{quote}

\subsection*{You'll want to know: \sk{प्रेरणा}}
\textbf{\sk{प्रेरणा}} — \emph{preraṇā} — \textquotedblleft{}inspiration; urging\textquotedblright{}.

\begin{quote}
\sk{यदि} \sk{गुरोः} \sk{प्रेरणा} \sk{अस्ति}, \sk{तर्हि} \sk{शिष्यः} \sk{उत्तमं} \sk{कार्यं} \sk{करोति।}
\end{quote}

\begin{quote}
\emph{yadi guroḥ preraṇā asti, tarhi śiṣyaḥ uttamaṃ kāryaṃ karoti.}
\end{quote}

\begin{quote}
With a teacher's urging, a pupil does fine work.
\end{quote}

A feminine noun.

\begin{cousinweb}
\textbf{\sk{प्रेरणा}} (\emph{preraṇā}) is built on \emph{pra}, \textquotedblleft{}forth\textquotedblright{}, and the root √īr, \textquotedblleft{}to set in motion\textquotedblright{}; MW gives \emph{preraṇa}, also \emph{preraṇā}, as setting in motion, urging, inciting; \textquotedblleft{}inspiration\textquotedblright{} is its modern use (MW).
\end{cousinweb}

\begin{grammarlens}[title={What you've built}]
One more word for saying what happens if.
\end{grammarlens}

\subsection*{Your turn}
\begin{itemize}
\raggedright
  \item \emph{Say it:} \emph{preraṇā}
  \item \emph{Say it:} \emph{preraṇā}, once more
  \item \emph{Say it:} \emph{yadi guroḥ preraṇā asti, tarhi śiṣyaḥ uttamaṃ kāryaṃ karoti.} — With a teacher's urging, a pupil does fine work.
  \item \emph{Recall:} say the Sanskrit for the time ahead, then the Sanskrit for a start, then say \emph{preraṇā} again
\end{itemize}

\begin{tcolorbox}[breakable,colback=teal!4,colframe=teal!35!black,title={Before you move on}]
What does \sk{प्रेरणा} mean? (\textquotedblleft{}Inspiration; urging\textquotedblright{}.) Say it once more.
\end{tcolorbox}
\section[\texorpdfstring{\sk{उद्युक्तः} (udyuktaḥ)}{उद्युक्तः (udyuktaḥ)}]{\sk{उद्युक्तः} (udyuktaḥ) — all set; prepared}

\label{lesson:SA-C555-udyuktah}



\begin{quote}
Before the new one: say the Sanskrit for a start, then the Sanskrit for inspiration.
\end{quote}

\subsection*{You'll want to know: \sk{उद्युक्तः}}
\textbf{\sk{उद्युक्तः}} — \emph{udyuktaḥ} — \textquotedblleft{}all set; prepared\textquotedblright{}.

\begin{quote}
\sk{यदि} \sk{सर्वे} \sk{उद्युक्ताः} \sk{सन्ति}, \sk{तर्हि} \sk{वयम्} \sk{अधुना} \sk{गच्छामः।}
\end{quote}

\begin{quote}
\emph{yadi sarve udyuktāḥ santi, tarhi vayam adhunā gacchāmaḥ.}
\end{quote}

\begin{quote}
If everyone is ready, we leave now.
\end{quote}

\begin{cousinweb}
\textbf{\sk{उद्युक्तः}} (\emph{udyuktaḥ}) is built on \emph{ud}, \textquotedblleft{}up\textquotedblright{}, and the root √yuj, \textquotedblleft{}to yoke\textquotedblright{}, as its past participle: yoked up; MW gives \emph{udyukta} as prepared or ready for, zealously active (MW).
\end{cousinweb}

\begin{grammarlens}[title={What you've built}]
One more word for saying what happens if.
\end{grammarlens}

\subsection*{Your turn}
\begin{itemize}
\raggedright
  \item \emph{Say it:} \emph{udyuktaḥ}
  \item \emph{Say it:} \emph{udyuktaḥ}, once more
  \item \emph{Say it:} \emph{yadi sarve udyuktāḥ santi, tarhi vayam adhunā gacchāmaḥ.} — If everyone is ready, we leave now.
  \item \emph{Recall:} say the Sanskrit for a start, then the Sanskrit for inspiration, then say \emph{udyuktaḥ} again
\end{itemize}

\begin{tcolorbox}[breakable,colback=teal!4,colframe=teal!35!black,title={Before you move on}]
What does \sk{उद्युक्तः} mean? (\textquotedblleft{}All set; prepared\textquotedblright{}.) Say it once more.
\end{tcolorbox}
\section[\texorpdfstring{\sk{प्रत्युत} (pratyuta)}{प्रत्युत (pratyuta)}]{\sk{प्रत्युत} (pratyuta) — on the contrary}

\label{lesson:SA-C555-pratyuta}



\begin{quote}
Before the new one: say the Sanskrit for inspiration, then the Sanskrit for all set.
\end{quote}

\subsection*{You'll want to know: \sk{प्रत्युत}}
\textbf{\sk{प्रत्युत}} — \emph{pratyuta} — \textquotedblleft{}on the contrary\textquotedblright{}.

\begin{quote}
\sk{यदि} \sk{वर्षा} \sk{भवति}, \sk{तर्हि} \sk{हानिः} \sk{न}, \sk{प्रत्युत} \sk{लाभः} \sk{भवति।}
\end{quote}

\begin{quote}
\emph{yadi varṣā bhavati, tarhi hāniḥ na, pratyuta lābhaḥ bhavati.}
\end{quote}

\begin{quote}
If it rains, it is no loss; on the contrary, it is a gain.
\end{quote}

\begin{cousinweb}
\textbf{\sk{प्रत्युत}} (\emph{pratyuta}) is built from \emph{prati}, \textquotedblleft{}back, against\textquotedblright{}, and \emph{uta}, \textquotedblleft{}and, even\textquotedblright{}; MW gives \emph{pratyuta} as on the contrary, rather (MW).
\end{cousinweb}

\begin{grammarlens}[title={What you've built}]
One more word for saying what happens if.
\end{grammarlens}

\subsection*{Your turn}
\begin{itemize}
\raggedright
  \item \emph{Say it:} \emph{pratyuta}
  \item \emph{Say it:} \emph{pratyuta}, once more
  \item \emph{Say it:} \emph{yadi varṣā bhavati, tarhi hāniḥ na, pratyuta lābhaḥ bhavati.} — If it rains, it is no loss; on the contrary, it is a gain.
  \item \emph{Recall:} say the Sanskrit for inspiration, then the Sanskrit for all set, then say \emph{pratyuta} again
\end{itemize}

\begin{tcolorbox}[breakable,colback=teal!4,colframe=teal!35!black,title={Before you move on}]
What does \sk{प्रत्युत} mean? (\textquotedblleft{}On the contrary\textquotedblright{}.) Say it once more.
\end{tcolorbox}
\section[\texorpdfstring{\sk{कर्मठः} (karmaṭhaḥ)}{कर्मठः (karmaṭhaḥ)}]{\sk{कर्मठः} (karmaṭhaḥ) — hard-working}

\label{lesson:SA-C555-karmathah}



\begin{quote}
Before the new one: say the Sanskrit for all set, then the Sanskrit for on the contrary.
\end{quote}

\subsection*{You'll want to know: \sk{कर्मठः}}
\textbf{\sk{कर्मठः}} — \emph{karmaṭhaḥ} — \textquotedblleft{}hard-working\textquotedblright{}.

\begin{quote}
\sk{यदि} \sk{त्वं} \sk{कर्मठः} \sk{असि}, \sk{तर्हि} \sk{तव} \sk{कार्यं} \sk{सम्यक्} \sk{भविष्यति।}
\end{quote}

\begin{quote}
\emph{yadi tvaṃ karmaṭhaḥ asi, tarhi tava kāryaṃ samyak bhaviṣyati.}
\end{quote}

\begin{quote}
If you are hard-working, your work will go well.
\end{quote}

\begin{cousinweb}
\textbf{\sk{कर्मठः}} (\emph{karmaṭhaḥ}) is built on \emph{karman}, \textquotedblleft{}work, action\textquotedblright{}, with the ending \emph{-ṭha}; MW gives \emph{karmaṭha} as capable of work, clever in work, working diligently (MW).
\end{cousinweb}

\begin{grammarlens}[title={What you've built}]
One more word for saying what happens if.
\end{grammarlens}

\subsection*{Your turn}
\begin{itemize}
\raggedright
  \item \emph{Say it:} \emph{karmaṭhaḥ}
  \item \emph{Say it:} \emph{karmaṭhaḥ}, once more
  \item \emph{Say it:} \emph{yadi tvaṃ karmaṭhaḥ asi, tarhi tava kāryaṃ samyak bhaviṣyati.} — If you are hard-working, your work will go well.
  \item \emph{Recall:} say the Sanskrit for all set, then the Sanskrit for on the contrary, then say \emph{karmaṭhaḥ} again
\end{itemize}

\begin{tcolorbox}[breakable,colback=teal!4,colframe=teal!35!black,title={Before you move on}]
What does \sk{कर्मठः} mean? (\textquotedblleft{}Hard-working\textquotedblright{}.) Say it once more.
\end{tcolorbox}
\section[\texorpdfstring{\sk{अध्यवस्यति} (adhyavasyati)}{अध्यवस्यति (adhyavasyati)}]{\sk{अध्यवस्यति} (adhyavasyati) — resolves}

\label{lesson:SA-C555-adhyavasyati}



\begin{quote}
Before the new one: say the Sanskrit for on the contrary, then the Sanskrit for hard-working.
\end{quote}

\subsection*{You'll want to know: \sk{अध्यवस्यति}}
\textbf{\sk{अध्यवस्यति}} — \emph{adhyavasyati} — \textquotedblleft{}resolves\textquotedblright{}.

\begin{quote}
\sk{यदि} \sk{सः} \sk{दृढम्} \sk{अध्यवस्यति}, \sk{तर्हि} \sk{सः} \sk{अवश्यं} \sk{जयति।}
\end{quote}

\begin{quote}
\emph{yadi saḥ dṛḍham adhyavasyati, tarhi saḥ avaśyaṃ jayati.}
\end{quote}

\begin{quote}
If he resolves firmly, he surely wins.
\end{quote}

Past (he …): \sk{सः} \sk{गन्तुम्} \sk{अध्यवसितवान्}, \emph{saḥ gantum adhyavasitavān}.

\begin{cousinweb}
\textbf{\sk{अध्यवस्यति}} (\emph{adhyavasyati}) is built on \emph{adhi}, \emph{ava} and the root √so, \textquotedblleft{}to bring to an end\textquotedblright{}; MW gives \emph{adhyavaso} as to undertake, to accomplish, to determine; its past participle is \emph{adhyavasita} (MW).
\end{cousinweb}

\begin{grammarlens}[title={What you've built}]
One more word for saying what happens if.
\end{grammarlens}

\subsection*{Your turn}
\begin{itemize}
\raggedright
  \item \emph{Say it:} \emph{adhyavasyati}
  \item \emph{Say it:} \emph{adhyavasyati}, once more
  \item \emph{Say it:} \emph{yadi saḥ dṛḍham adhyavasyati, tarhi saḥ avaśyaṃ jayati.} — If he resolves firmly, he surely wins.
  \item \emph{Recall:} say the Sanskrit for on the contrary, then the Sanskrit for hard-working, then say \emph{adhyavasyati} again
\end{itemize}

\begin{tcolorbox}[breakable,colback=teal!4,colframe=teal!35!black,title={Before you move on}]
What does \sk{अध्यवस्यति} mean? (\textquotedblleft{}Resolves\textquotedblright{}.) Say it once more.
\end{tcolorbox}
\section[(dialogue)]{Looking back: from Longings to The Road Ahead}

\label{lesson:SA-R555-review-549-555}



\begin{quote}
Say the Sanskrit for is eager for and a craving.
\end{quote}

\subsection*{Your turn}
Say these aloud:

\begin{itemize}
\raggedright
  \item is eager for, a craving, a wait, a wish to see, strives for
  \item divides up, each one, an inheritance, puts to use, a handing out
  \item leaves behind, regrets, remorse, detached, lets go of
  \item practises regularly, a daily routine, exertion, makes use of, helpful
  \item a slip, come now, guards well, a reversal, a reproach
  \item is possible, far-sighted, intends, the time ahead, a start
  \item inspiration, all set, on the contrary, hard-working, resolves
\end{itemize}

\begin{tcolorbox}[breakable,colback=teal!4,colframe=teal!35!black,title={Before you move on}]
Is eager for? (\textbf{\sk{स्पृहयति}}, \emph{spṛhayati}.) A craving? (\textbf{\sk{लालसा}}, \emph{lālasā}.) A wait? (\textbf{\sk{प्रतीक्षा}}, \emph{pratīkṣā}.) A wish to see? (\textbf{\sk{दिदृक्षा}}, \emph{didṛkṣā}.) Strives for? (\textbf{\sk{समीहते}}, \emph{samīhate}.) Divides up? (\textbf{\sk{वण्टयति}}, \emph{vaṇṭayati}.) Each one? (\textbf{\sk{प्रत्येकम्}}, \emph{pratyekam}.) An inheritance? (\textbf{\sk{दायः}}, \emph{dāyaḥ}.) Puts to use? (\textbf{\sk{विनियुंक्ते}}, \emph{viniyuṅkte}.) A handing out? (\textbf{\sk{वितरणम्}}, \emph{vitaraṇam}.) Leaves behind? (\textbf{\sk{जहाति}}, \emph{jahāti}.) Regrets? (\textbf{\sk{अनुशोचति}}, \emph{anuśocati}.) Remorse? (\textbf{\sk{पश्चात्तापः}}, \emph{paścāttāpaḥ}.) Detached? (\textbf{\sk{विरक्तः}}, \emph{viraktaḥ}.) Lets go of? (\textbf{\sk{उत्सृजति}}, \emph{utsṛjati}.) Practises regularly? (\textbf{\sk{परिशीलयति}}, \emph{pariśīlayati}.) A daily routine? (\textbf{\sk{दिनचर्या}}, \emph{dinacaryā}.) Exertion? (\textbf{\sk{प्रयासः}}, \emph{prayāsaḥ}.) Makes use of? (\textbf{\sk{प्रयुंक्ते}}, \emph{prayuṅkte}.) Helpful? (\textbf{\sk{उपकारकः}}, \emph{upakārakaḥ}.) A slip? (\textbf{\sk{स्खलनम्}}, \emph{skhalanam}.) Come now? (\textbf{\sk{हन्त}}, \emph{hanta}.) Guards well? (\textbf{\sk{परिरक्षति}}, \emph{parirakṣati}.) A reversal? (\textbf{\sk{विपर्ययः}}, \emph{viparyayaḥ}.) A reproach? (\textbf{\sk{उपालम्भः}}, \emph{upālambhaḥ}.) Is possible? (\textbf{\sk{सम्भवति}}, \emph{sambhavati}.) Far-sighted? (\textbf{\sk{दीर्घदर्शी}}, \emph{dīrghadarśī}.) Intends? (\textbf{\sk{अभिप्रैति}}, \emph{abhipraiti}.) The time ahead? (\textbf{\sk{आयतिः}}, \emph{āyatiḥ}.) A start? (\textbf{\sk{आरम्भः}}, \emph{ārambhaḥ}.) Inspiration? (\textbf{\sk{प्रेरणा}}, \emph{preraṇā}.) All set? (\textbf{\sk{उद्युक्तः}}, \emph{udyuktaḥ}.) On the contrary? (\textbf{\sk{प्रत्युत}}, \emph{pratyuta}.) Hard-working? (\textbf{\sk{कर्मठः}}, \emph{karmaṭhaḥ}.) Resolves? (\textbf{\sk{अध्यवस्यति}}, \emph{adhyavasyati}.)
\end{tcolorbox}
\section[(dialogue)]{The Road Ahead — once more}

\label{lesson:SA-R555-review-final-555}



\begin{quote}
Say the Sanskrit for inspiration and all set.
\end{quote}

\subsection*{Your turn}
\begin{itemize}
\raggedright
  \item \emph{Say it:} inspiration, all set, on the contrary, hard-working, resolves
\end{itemize}

\begin{tcolorbox}[breakable,colback=teal!4,colframe=teal!35!black,title={Before you move on}]
Inspiration? (\textbf{\sk{प्रेरणा}}, \emph{preraṇā}.) All set? (\textbf{\sk{उद्युक्तः}}, \emph{udyuktaḥ}.) On the contrary? (\textbf{\sk{प्रत्युत}}, \emph{pratyuta}.) Hard-working? (\textbf{\sk{कर्मठः}}, \emph{karmaṭhaḥ}.) Resolves? (\textbf{\sk{अध्यवस्यति}}, \emph{adhyavasyati}.)
\end{tcolorbox}
